% ===== ≈ 0,5 trang =====
\begin{table}[H]
\centering
\caption{Các bộ dữ liệu sử dụng.}\label{tab:datasets}
\small
\begin{tabularx}{\linewidth}{clrX}
\toprule
Bộ & Nguồn & Ảnh gốc & Dùng cho \\
\midrule
1 & VAIPE (đơn thuốc Việt Nam, Kaggle) & \num{1345} & Tách khoản thuốc; lớp đơn thuốc; nguồn ảnh sạch của \bo{5} \\
2 & Roboflow Report Keywords Detection (CC BY 4.0) & 57 & Phiếu xét nghiệm (tiếng Anh): dựng bảng; lớp phiếu XN \\
3 & RVL-CDIP~\cite{harley2015rvlcdip} (tập con) & \num{1000} & Lớp ``tài liệu khác'' (800) + lớp lạ open-set (200) \\
4 & Nhóm tự chụp & 83 & Chỉ để test: phân loại (định lượng), các bước khác (định tính) \\
5 & Tổng hợp từ ảnh sạch \bo{1} & \num{1500} & Ground truth chính xác cho mọi bước \\
\bottomrule
\end{tabularx}
\end{table}
Nhãn có sẵn: \bo{1} có khung và chữ từng dòng với vai trò \texttt{drugname}, \texttt{usage}, \texttt{quantity},~\ldots;
\bo{2} có khung từng cụm từ với vai trò (tên chỉ số, kết quả, đơn vị, khoảng tham chiếu), ảnh đã bị ép về $640\times640$;
\bo{4} gồm 48 phiếu xét nghiệm, 18 đơn thuốc, 17 tài liệu khác, chỉ có nhãn loại tài liệu.
Ground truth khoản thuốc của \bo{1} được gom tự động từ nhãn dòng: mỗi \texttt{drugname} cùng các box \texttt{usage},
\texttt{quantity} cùng hàng hoặc ngay dưới; khoản là dải ngang trải hết bề ngang vùng thuốc (nhãn đôi khi thiếu box số lượng).
Ground truth dòng / cột của \bo{2} suy từ nhãn cụm từ: hai cụm từ cùng dòng nếu chồng nhau theo chiều dọc $\ge 50\%$; cột
là vai trò cụm từ. Sau khi soát bằng mắt cả 55 ảnh có nhãn thân bảng, giữ \textbf{32 ảnh} đủ cơ sở (loại ảnh nhãn lệch khỏi chữ,
ảnh mờ không đọc được, bố cục máy in nhiều cột); 9 ảnh được sửa dòng theo quy tắc ghi trong \texttt{k0\_review.json} (bỏ dòng
mô tả và tiêu đề nhóm, gộp tên chỉ số hai dòng). \Bo{2} không có nhãn chữ: chữ từng khung do LLM phiên âm từ ảnh gốc rồi
người soát (mục~7), dùng làm đáp án CER cho K5.
